\documentclass[spanish,letterpaper,oneside]{article}

\def\documenttitle {Universo 3-60 · Ficha de trabajo}
\def\documentsubtitle {Temas selectos de matemáticas I · Cohorte SD26}
\def\documentsubject {Actividad 3 · Material para estudiantes}
\def\documentauthor {Mtro. Martín Jonathan de la Cruz Muñoz}
\def\coursename {Temas selectos de matemáticas I}
\def\coursecode {PFM02}
\def\universityname {Instituto de Investigación, Innovación y Estudios de Posgrado para la Educación del Estado de Nuevo León}
\def\universityfaculty {Maestría en Aprendizaje y Enseñanza de las Matemáticas}
\def\universitydepartment {Temas selectos de matemáticas I}
\def\universitydepartmentimage {img/iiiepe/marca-iiiepe.png}
\def\universitydepartmentimagecfg {height=1.75cm}
\def\universitylocation {Monterrey, Nuevo León}
\def\studentfullname {Mtro. Martín Jonathan de la Cruz Muñoz}
\def\studentid {00874}
\def\studentemail {martin.delacruz@campus.iiiepe.edu.mx}
\def\programperiod {2do periodo, agosto--noviembre de 2026 (SD26)}
\def\programmodality {Modalidad presencial y en línea}
\def\programshort {MAEM}

\def\authortable {
  \begin{tabular}{ll}
    \textbf{Docente:} & \studentfullname \\
    Matrícula: & \studentid \\
    Unidad: & \coursename \\
    Periodo: & \programperiod \\
    \multicolumn{2}{l}{Fecha: 24 de octubre de 2026}
  \end{tabular}
}

\input{template}

\usepackage{helvet}
\renewcommand{\familydefault}{\sfdefault}
\usepackage{amsmath,amssymb,array,longtable,tikz,tabularx,colortbl,enumitem}
\definecolor{iiiepeNavy}{HTML}{163B5C}
\definecolor{iiiepeBlue}{HTML}{1F6F9F}
\definecolor{iiiepePale}{HTML}{EEF5F8}
\def\portraittitlecolor {iiiepeNavy}

\setlength{\parskip}{5pt}\setlength{\parindent}{0pt}
\begin{document}\templatePortrait\templatePagecfg\templateFinalcfg\fancyfoot{}\fancyfoot[L]{\scriptsize IIIEPE · MAEM · SD26}\fancyfoot[C]{\scriptsize\thepage}\fancyfoot[R]{\scriptsize Ficha Act. 3}\setstretch{1.2}\raggedbottom

\section*{Triangulación de polígonos y ángulos en la circunferencia}
Nombre: \rule{8cm}{.2pt}\hfill Fecha: \rule{2.5cm}{.2pt}

\subsection*{1. Regla general de suma de ángulos interiores}
Completa la tabla trazando las diagonales desde un solo vértice y deduce la expresión para $n$ lados:

\begin{center}
\begin{tabular}{|l|c|c|c|}\hline
\textbf{Polígono} & \textbf{Lados ($n$)} & \textbf{Triángulos} & \textbf{Suma interior} \\\hline
Hexágono & $6$ & & \\[0.45cm]\hline
Octágono & $8$ & & \\[0.45cm]\hline
Decágono & $10$ & & \\[0.45cm]\hline
Icoságono & $20$ & & \\[0.45cm]\hline
\end{tabular}
\end{center}

\subsection*{2. Descubrimiento en GeoGebra}
Construye un ángulo central $\angle AOB$ y un ángulo inscrito $\angle ACB$ sobre el mismo arco. Mueve $C$ sin pasar por $A$ ni $B$ y registra qué relación exacta cumplen sus medidas. Justifica con un esquema usando triángulos isósceles de radio $r$.
\vspace{2.2cm}

\end{document}
